\scanpage{28}
% lines 1–6
à la donnée d'un fibré en ``plans affines symplectiques'' (ou unimodulaires). Le \inserted{\unsure{schém. en|schéma en}} groupes des autom. d'une telle structure est donc le produit $1/2$ direct de $\underline{\mathrm{SL}}(V)$ et de $\underline V$ -- lisse de dim 5 -- ce n'est en tous cas pas \deleted{\unsure{le}} une forme de $\mathrm{GP}(1)$ !

% lines 6–9
Le groupe des autom. \deleted{\unsure{de}}d'une forme de $\mathfrak{sl}(2)$ ou $\mathfrak{gp}(1)$ est $\simeq$ groupe affine à deux variables, lisse de dim rel. 6.

% lines 10–12
On déduit aisément, de la structure précédente, les données $f,f',\rho_{E'}$. En effet, $f$ est le composé

% lines 12
\diagram{\begin{tikzcd}[ampersand replacement=\&,column sep=small]
E \arrow[r,"\mathrm{can}"] \& V \arrow[r,"\varphi"] \& V' \arrow[r,"\mathrm{can}"] \& E',
\end{tikzcd}}

% lines 12–13
$f'$ est le composé

% lines 13
\diagram{\begin{tikzcd}[ampersand replacement=\&,column sep=small]
E' \arrow[r,"\mathrm{can}"] \& \mathscr O_S \arrow[r,no head,"\varphi" description] \& \det V \arrow[r,"\mathrm{can}"] \& E,
\end{tikzcd}}

% lines 13–15
enfin $\rho_{E'}$ \inserted{(op. de $E'$ sur $E$)} est l'opération \deleted{qui sur $V'=f(E)$ se} \inserted{contragrédiente de l'op. adjointe de $E'$} \ldots.

% lines 15–19
\deleted{\unsure{déduit} de la représentation adjointe de $E$, qui est nulle sur $\det V=\mathscr O_S\subset E$, et qui pour un $H\in E'$ au dessus de 1 donne l'opération $\rho_{E'}(H):E\mathbin{-}E$ qui (étant nulle sur $\det V=\mathscr O_S$) se factorise en \deleted{\unsure{le}}\inserted{un} splitting}

% lines foot
% archive sheet number: 27
